% problem_id: S051-drops-D2
% source_id: S051 (DROPS/LIPIcs)
% source_file: t5.pdf
% statement_locator: printed p. 11
% proof_locator: printed p. [11, 12, 13]
% representation: PDF; transcribed from rendered page images (never from the text layer)
% collected_by: carrier rule (the headline theorem was an assembly)
% review_status: two independent reviewers, H2 agreed
% status: complete (statement + author proof verbatim + reason + locators)
%
% =====================================================================
%  D2  --  verbatim transcription of ONE complete theorem-proof unit
% ---------------------------------------------------------------------
%  Source PDF : /disks/sata1/yupeng/human-proof-corpus/source-cache/
%               registry-sources/drops/t5.pdf   (20 pages, host Yupeng-orx)
%  Paper      : F. Abbasi, M. Adamczyk, M. Bosch-Calvo, J. Byrka,
%               F. Grandoni, K. Sornat, A. Tinguely,
%               "An O(log log n)-Approximation for Submodular Facility
%               Location", LIPIcs Vol. 297, ICALP 2024, Article No. 5,
%               pp. 5:1--5:20.  DOI 10.4230/LIPIcs.ICALP.2024.5
%  UNIT       : Lemma 3.3 (CP-rounding for DLA) TOGETHER WITH its complete
%               proof, which in the paper consists of Algorithm 1 +
%               Lemma 3.5 (feasibility) + Lemma 3.6 (approximation ratio).
%               Printed pages 5:11, 5:12, 5:13  =  PDF pages 11, 12, 13.
%  METHOD     : pdftotext -layout was used ONLY to locate items by page.
%               Every symbol below was read from pdftoppm renders
%               (300 dpi page bands, 600 dpi tight crops) via read_image.
%               The text layer is unreliable here: it mangles \tilde F /
%               F, \hat h / h, \Delta_i^\theta, and it renders the small
%               caps / blackboard glyphs inconsistently.
%  VERBATIM, no rewriting, no completion.  Original typos are copied and
%  flagged with %% [NOTE: ...] / %% [TYPO? ...] comments.  Nothing was
%  silently corrected.
% =====================================================================
\documentclass[11pt]{article}
\usepackage[a4paper,margin=2.4cm]{geometry}
\usepackage{amsmath,amssymb}
\usepackage[T1]{fontenc}
\usepackage{lmodern}

\newcommand{\cost}{\operatorname{cost}}
\newcommand{\costDLA}{\operatorname{cost}_{\mathrm{DLA}}}
\newcommand{\conn}{\operatorname{conn}}
\newcommand{\openop}{\operatorname{open}}
\newcommand{\DLA}{\mathrm{DLA}}
\newcommand{\tF}{\tilde F}
\newcommand{\tC}{\tilde C}
\newcommand{\tT}{\tilde T}
\newcommand{\tphi}{\tilde\varphi}
\newcommand{\hatt}[1]{\hat{#1}}

% paper-style headline item:  blacktriangleright  Lemma X.  <statement>
\newcommand{\pitem}[2]{\par\medskip\noindent$\blacktriangleright$\,\textbf{#1}\;\textit{#2}\par\medskip}

\begin{document}

\begin{center}
{\large\bfseries Lemma 3.3 of Abbasi--Adamczyk--Bosch-Calvo--Byrka--Grandoni--Sornat--Tinguely,\\
``An O$(\log\log n)$-Approximation for Submodular Facility Location'', ICALP 2024}
\end{center}

\medskip
\noindent\rule{\textwidth}{0.6pt}

% ---------------------------------------------------------------------
%  CONTEXT (NOT part of the unit; quoted only to show what Lemma 3.3 is
%  the carrier of).  Each item below was read at 600 dpi from its own page.
% ---------------------------------------------------------------------
\section*{Context read off other pages (not part of the transcribed unit)}

\noindent\emph{Theorem 1.1, printed page 5:3 = PDF page 3:}
\pitem{Theorem 1.1.}{There is a polynomial-time $O(\log\log n)$-approximation algorithm for SFL.}

\noindent\emph{Lemma 3.1, printed page 5:9 = PDF page 9:}
\pitem{Lemma 3.1.}{Given a feasible fractional solution $x$ to (Conf-LP) for an HST-type SFL
instance, in polynomial time one can compute a feasible (integral) solution for the same
instance with cost at most $O(\log\log N)\cdot\cost(x)$.}

\noindent\emph{Lemma 3.2, printed page 5:10 = PDF page 10:}
\pitem{Lemma 3.2.}{Given a polynomial-time $O(\log D)$-approximate CP-rounding algorithm for
DLA w.r.t.\ (DLA-CP), where $D$ is the depth of the tree, there is polynomial-time
$O(\log\log N)$-approximate LP-rounding algorithm for SFL on HST-type instances with
tree-depth $O(\log N)$ w.r.t.\ (Conf-LP).}
%% [NOTE: the article "a" is missing before "polynomial-time
%%  O(log log N)-approximate LP-rounding algorithm" in the original.
%%  Copied verbatim, not corrected.]
%% [CONTEXT, in my own words, NOT a transcription: DLA (Descendant-Leaf
%%  Assignment) is defined in Section 3.1 (printed page 5:9): a rooted tree
%%  \tilde T of depth D, facilities \tilde F and clients \tilde C each mapped
%%  to a node, facilities to leaves; \tilde F_c = facilities mapped below v(c);
%%  a solution assigns each c to f in \tilde F_c and costs
%%  sum_f h(\tilde\varphi^{-1}(f)) with h monotone submodular, h(empty)=0.
%%  (DLA-CP) is min sum_f \hat h(z^f) s.t. sum_{f in \tilde F_c} z_c^f = 1,
%%  z_c^f >= 0.]

\noindent\emph{Proof of Theorem 1.1, printed page 5:9 = PDF page 9} chains, with no
computation of its own: Lemma 1.7, Lemma 1.8, the construction of Section 2, Lemma 2.2,
the FRT embedding (Theorem 1.9) and finally Lemma 3.1. Lemma 3.1 in turn ``follows by
chaining Lemmas 3.2 and 3.3'' (printed page 5:11). Lemma 3.2 is \emph{conditional}: it
\emph{assumes} an $O(\log D)$-approximate CP-rounding algorithm for DLA. Lemma 3.3 is the
node at which that assumption is discharged, i.e.\ the carrier collected below.

\noindent\rule{\textwidth}{0.6pt}

% =====================================================================
%  THE UNIT
% =====================================================================
\section*{Transcription}

%% ===================== page 11  (printed 5:11) =====================
%% --- page 11 ---
\noindent\textit{\small [running head: F.~Abbasi et al.\hfill 5:11]}

\medskip
%% [CONTEXT ONLY: the paragraph below is the last paragraph of the proof of
%%  Lemma 3.2; it is transcribed because it stands at the top of the same
%%  page, but it does NOT belong to the collected unit.]
Therefore, the connection cost of $c$ in $\varphi$ is at most $2$ times its connection cost in $x$. We
conclude that $\conn(\varphi)\le 2\conn(x)$. Altogether we have $\cost(\varphi)\le 2\conn(x)+O(\log\log N)\cdot
2\openop(x)\le O(\log\log N)\cdot\cost(x)$.\hfill$\blacktriangleleft$
%% [NOTE: the two occurrences of "O(log log N)" on this line were checked at
%%  600 dpi; an initial 150 dpi read of the same line gave "O(log N)".]

\subsection*{3.2\quad An Approximation Algorithm for DLA}

In this section, we present a CP-rounding algorithm for DLA. Lemma 3.1 follows by chaining
Lemmas 3.2 and 3.3.

\pitem{Lemma 3.3.}{Given a feasible fractional solution $z$ to (DLA-CP) on an instance of DLA
with tree-depth $D$, in polynomial time one can compute a feasible (integral) solution to the
same instance of cost at most $O(\log D)\cdot\costDLA(z)$.}

The CP-rounding algorithm from Lemma 3.3 is essentially the algorithm by Bosman and
Olver [5] with minor modifications that we introduced to simplify our correctness analysis.
Also, the analysis of its approximation ratio is essentially identical to [5], but we reproduce it
for the sake of completeness. In particular, we will exploit the following definitions and lemma
from [5]. Let $h:2^{\tC}\to\mathbb{R}_{\ge 0}$ be a monotone submodular function with $h(\emptyset)=0$. For a given
$f\in\tF$ and a (possibly infeasible) solution $z$ to (DLA-CP), let $L_\theta(z^f):=\{c\in\tC:z_c^f\ge\theta\}$ be
the set of clients that are served fractionally by at least some value $\theta$ by $f$. Let also $z^{f|\theta}$
be obtained from $z^f$ by rounding down to $\theta$ the values larger than $\theta$, i.e.\ $z_c^{f|\theta}:=\min\{z_c^f,\theta\}$
for each $c\in\tC$. Given $\theta\in[0,1]$ and $z^f\in[0,1]^{\tC}$, we say that the set $L_\theta(z^f)$ is $\alpha$-supported
(w.r.t.\ $h$) if $\hatt h(z^f)-\hatt h(z^{f|\theta})\ge\alpha h(L_\theta(z^f))$.

\pitem{Lemma 3.4 (Lemma 5.2 from [5]).}{Given $z^f\in[0,1]^{\tC}$ and $\alpha\in(0,1]$, at least one of the
following holds: (1) there exists $\theta\in[0,1]$, which can be computed in polynomial time, such
that $L_\theta(z^f)$ is $\frac{\alpha}{32}$-supported; (2) $2^{1/\alpha}h(L_1(z^f))\le\hatt h(z^f)$.}

Our algorithm is Algorithm 1 in the figure. Recall that $\tT_v$ is the subtree rooted at node $v$,
where $\tT_v$ includes $v$ and all its descendants. Furthermore, $\tF_v$ is the set of facilities mapped
to the leaves of $\tT_v$. As usual the level of a node is its hop-distance from the root.

\bigskip
\noindent\rule{\textwidth}{0.4pt}
\noindent\textbf{Algorithm 1} An algorithm used to prove Lemma 3.3.
\noindent\rule{\textwidth}{0.4pt}
\begin{quote}
\textbf{Input:} Feasible solution $z$ to (DLA-CP)\\
1:\ $S^f\leftarrow\emptyset$ for all $f\in F$\\
%% [TYPO? Algorithm 1 line 1 prints plain "F", while the facility set of a DLA
%%  instance is \tilde F everywhere else on this page and on pages 10 and 12.
%%  Copied verbatim as printed; not corrected.]
2:\ \textbf{for} $i=0,\dots,D$ \textbf{do}\\
3:\quad For every node $v$ at level $D-i$, choose an arbitrary $f_v\in\tF_v$ and set $z^{f_v}\leftarrow\sum_{f'\in\tF_v}z^{f'}$ and
$z^{f'}\leftarrow 0$ for all $f'\in\tF_v\setminus\{f_v\}$\\
4:\quad \textbf{if} there exists $\theta\in[0,1]$ such that $L_\theta(z^{f_v})$ is $\frac{1}{32\log(D+1)}$-supported \textbf{then}\\
5:\quad\quad For an arbitrary such $\theta$, set $S^{f_v}\leftarrow S^{f_v}\cup L_\theta(z^{f_v})$ and $z_c^{f_v}\leftarrow 0$ for all $c\in L_\theta(z^{f_v})$\\
6:\quad \textbf{else}\\
7:\quad\quad Set $S^{f_v}\leftarrow S^{f_v}\cup L_1(z^{f_v})$ and $z_c^{f_v}\leftarrow 0$ for all $c\in L_1(z^{f_v})$\\
8:\ For every $c\in\tC$ choose $f\in\tF_c$ such that $c\in S^f$ and set $S^{f'}\leftarrow S^{f'}\setminus\{c\}$ for all $f'\in\tF\setminus\{f\}$\\
9:\ \textbf{return} $(S^f)_{f\in\tF}$
\end{quote}
\noindent\rule{\textwidth}{0.4pt}
\bigskip

Clearly Algorithm 1 runs in polynomial time. The next two lemmas analyze the correctness
and the approximation ratio of Algorithm 1, hence proving Lemma 3.3.

\pitem{Lemma 3.5.}{Algorithm 1 computes a feasible DLA solution.}

\noindent\textbf{Proof.}\ Consider a given client $c\in\tC$ such that $v(c)$ is at level $D-i$ in $\tT$. Let us show that
the following invariant holds at the beginning of each iteration $j\le i$: either $\sum_{f\in\tF_c}z_c^f=1$
or $c\in S^f$ for some $f\in\tF_c$. The invariant trivially holds for $j=0$. Assume that it holds

%% ===================== page 12  (printed 5:12) =====================
%% --- page 12 ---
\noindent\textit{\small [running head: 5:12\hfill An $O(\log\log n)$-Approximation for Submodular Facility Location]}

\medskip
up to the beginning of iteration $j<i$, and consider what happens during that iteration.
Notice that for every node $v$ at level $D-j>D-i$, we either have that every $f\in\tF_v$ is a
descendant of $v(c)$ or every $f\in\tF_v$ is not in $\tF_c$. Therefore, in Step (3) the value of $\sum_{f\in\tF_c}z_c^f$
does not change. In more detail, it remains $1$ by inductive hypothesis. The same value can
decrease in Steps (5) or (7), however, this can only happen if $c$ is added to $S^{f_v}$ for some
$f_v\in\tF_c$. Thus the invariant holds at the end of the $j$-th iteration, hence at the beginning of
the next iteration $j+1$.

Due to the invariant, during the iteration $i$, when one considers the node $v=v(c)$, one
has that either $c$ already belongs to some $S^f$ with $f\in\tF_c$, or $\sum_{f\in\tF_c}z_c^f=1$. In the latter
case, after Step (3), $z_c^{f_v}=1$ where $f_v\in\tF_c$, so $c$ belongs to every set $L_\theta(z^{f_v})$ with $\theta\in[0,1]$.
As a consequence, $c$ is added to $S^{f_v}$ either in Step (5) or in Step (7).

It might happen that a client $c$ is assigned \textit{also} to a facility not in $\tF_c$. Step (8) guarantees
that the final assignment of $c$ is correct and unique.\hfill$\blacktriangleleft$

\pitem{Lemma 3.6.}{Algorithm 1 outputs a solution of cost at most $O(\log D)\cdot\costDLA(z)$.}

\noindent\textbf{Proof.}\ Recall that $\costDLA(z)=\sum_{f\in\tF}\hatt h(z^f)$. We start by observing that the value of
$\costDLA(z)$ can not increase over time when $z$ changes during the execution of the algorithm.
Indeed, Steps (5) and (7) can only decrease the entries of $z$, hence $\costDLA(z)$ by the
monotonicity of $\hatt h(\cdot)$. The only other changes of $z$ happen in Step (3). Let us interpret
this step as iteratively decreasing to zero $z^{f'}$ for each $f'\in\tF_v\setminus\{f_v\}$ and increasing $z^{f_v}$ by
the same amount. The decrease of the cost at each step is $\hatt h(z^{f_v})+\hatt h(z^{f'})-\hatt h(z^{f_v}+z^{f'})$.
By the alternative definition of $\hatt h(\cdot)$ as in (2) and its convexity, one has $\hatt h(z^{f_v}+z^{f'})=
2\hatt h\!\left(\frac{z^{f_v}+z^{f'}}{2}\right)\le 2\left(\frac12\hatt h(z^{f_v})+\frac12\hatt h(z^{f'})\right)=\hatt h(z^{f_v})+\hatt h(z^{f'})$. Hence the decrease of the cost is
non-negative as required.

For each facility $f$ and level $i$, let $\Delta_i^\theta(f)$ be the clients added to $S^f$ in Step (5) during
iteration $i$ (possibly $\Delta_i^\theta(f)=\emptyset$). We define similarly $\Delta_i^1(f)$ w.r.t.\ Step (7). Notice that, by
the submodularity (hence subadditivity) of $h(\cdot)$, the increase of the cost of the solution due
to adding $\Delta$ to $S^f$ is at most $h(\Delta)$. Therefore we can upper bound the cost of the final
solution $S=(S^f)_{f\in\tF}$ by
\[
\costDLA(S):=\sum_{f\in\tF}h(S^f)\;\le\;\sum_{i=0}^{D}\sum_{f\in\tF}\Big(h(\Delta_i^\theta(f))+h(\Delta_i^1(f))\Big).
\]

Let us upper bound the right-hand side of the above inequality. Let $z(i)$ denote the value of
$z$ at the beginning of iteration $i$. From the previous observation, we have $\hatt h(z(i))\le\hatt h(z)$ for
every $i$. By Lemma 3.4 with $\alpha=\frac{1}{\log(D+1)}$, for any $\Delta_i^1(f)$ one has $h(\Delta_i^1(f))\le\frac{1}{D+1}\hatt h(z^f(i))$.
Thus
\begin{equation}
\sum_{i=0}^{D}\sum_{f\in\tF}h(\Delta_i^1(f))\;\le\;\sum_{i=0}^{D}\sum_{f\in\tF}\frac{1}{D+1}\hatt h(z^f(i))\;\le\;\sum_{i=0}^{D}\frac{1}{D+1}\costDLA(z(i))\;\le\;\costDLA(z).\tag{4}
\end{equation}

Let $z(D+1)$ be the value of $z$ at the end of the $D$-th iteration, hence in particular
$\costDLA(z(D+1))\ge 0$. Notice that $z=z(0)$. We can lower bound $\costDLA(z)$ by
\[
\costDLA(z)\;\ge\;\sum_{i=0}^{D}\Big(\costDLA(z(i))-\costDLA(z(i+1))\Big).
\]

%% ===================== page 13  (printed 5:13) =====================
%% --- page 13 ---
\noindent\textit{\small [running head: F.~Abbasi et al.\hfill 5:13]}

\medskip
Let $z_1(i)$ be the value of $z$ obtained from $z(i)$ after applying Step (3) for all nodes of level
$D-i$. Let also $z_2(i)$ be the value obtained from $z_1(i)$ if, for all the facilities $F_i'$ where
Step (5) is applied during iteration $i$, instead of setting $z_c^f=0$ one sets $z_c^f=\theta$ for the
corresponding value of $\theta$. For the facilities not in $F_i'$ we simply let $z_2^f(i)=z_1^f(i)$. Observe
that $z(i+1)\le z_2(i)\le z_1(i)\le z(i)$. One has
\begin{align*}
\costDLA(z(i))-\costDLA(z(i+1)) &\ge \costDLA(z_1(i))-\costDLA(z(i+1))\\
&\ge \costDLA(z_1(i))-\costDLA(z_2(i))\\
&=\sum_{f\in\tF}\hatt h\!\left(z_1^f(i)\right)-\hatt h\!\left(z_2^f(i)\right)=\sum_{f\in F_i'}\hatt h\!\left(z_1^f(i)\right)-\hatt h\!\left(z_2^f(i)\right)\\
&\ge\frac{\sum_{f\in F_i'}h\!\left(\Delta_i^\theta(f)\right)}{32\log(D+1)}=\frac{\sum_{f\in\tF}h\!\left(\Delta_i^\theta(f)\right)}{32\log(D+1)}.
\end{align*}

In the first two inequalities above we used the monotonicity of $\hatt h(\cdot)$, while in the last inequality
the definition of $\alpha$-supported. Altogether
\begin{equation}
\sum_{i=0}^{D}\sum_{f\in\tF}h\!\left(\Delta_i^\theta(f)\right)\;\le\;32\log(D+1)\cdot\sum_{i=0}^{D}\Big(\costDLA(z(i))-\costDLA(z(i+1))\Big)
\;\le\;O(\log D)\cdot\costDLA(z).\tag{5}
\end{equation}

By the monotonicity of $h(\cdot)$, Step (8) cannot increase the cost of the solution, hence the
claim.\hfill$\blacktriangleleft$

\medskip
\noindent\textit{[End of the collected unit.  The next line of the paper, still on printed page
5:13 = PDF page 13, is the section heading ``4\quad Universal Stochastic Facility Location''.]}

\end{document}
